\begin{statement}{A.30}{Proved}{The explicit Collino unit removes the stable genus-three exception}
Let $\ell$ be odd. For every original packet whose action is
induced on the stable $S_3$ cover
\[
                  C^\circ:\quad v^2=w^7-w,\qquad w\ne0,\infty,
\]
and every permitted open subgroup $\Delta_0$, one has
\[
                    P_{C,\ell}^{\,s_b(\Delta_0)}=I_{C,\ell}.
                                                        \tag{A.30.1}
\]
In particular the possible zero-first-Johnson packet in A.29
does not exist. This conclusion includes $\ell=3$.

The arithmetic input is an explicit higher-Chow pushforward.
Put $L=\mathbb Q(i,\sqrt3)$, $\zeta=(-1+i\sqrt3)/2$,
$\epsilon=2+\sqrt3$, and $F:Y^2=X^3+3X$. There is a
degree-four map $q:C\to F$ for which the Collino cycle of
the two removed Weierstrass points pushes forward to
\[
                  \zeta^2\epsilon^{-4}\in
                    \operatorname{CH}^1(F,1)=L^\times.
                                                        \tag{A.30.2}
\]
After subtracting twice the degree-two reflection
pushforward, this is a functional on the primitive
exterior-square representation. Its rational
$\ell$-adic regulator is $-4\kappa_\epsilon$, which
remains nonzero on every open subgroup of $\Delta$.
The comparison with the actual Johnson class is proved
below on universal moduli and then specialized; it does
not assume faithfulness of a Hodge realization.
\end{statement}
\begin{proof}
\paragraph{1. A canonical universal extension for the twice-punctured curve.}
We first work on the moduli stack $\mathcal H$ of smooth
hyperelliptic curves of genus $g\ge2$ with all
Weierstrass points labelled. Choose three of its
sections, denoted by $P,Q,R$. Write $\mathcal C^\circ
=\mathcal C\setminus(Q\cup R)$, and use the ordinary
section $P$ as base point. The quotient by the
hyperelliptic involution is a projective-line bundle
with three disjoint sections, so it has the unique
coordinate $f$ with values $0,\infty,1$ at $Q,R,P$.
On each fibre,
\[
                         \operatorname{div}(f)=2Q-2R.
\]
All these constructions are defined over $\mathbb Q$.
They may equally be performed on a finite level cover
and descended; the stack notation keeps the geometric
mapping class group unchanged.

Let $\mathcal V=H_1(\mathcal C/\mathcal H,\mathbb Q_\ell)$,
let $t$ be the positive local class about $Q$, and
let $\mathfrak p$ be the relative rational unipotent
Lie algebra of $\mathcal C^\circ$ based at $P$.
Use the ordinary lower central series $\Gamma_\bullet$.
The hyperelliptic involution $h$ fixes $P$ and is an
actual involution commuting with all monodromy.
Its eigenspaces on $H_1(\mathcal C^\circ)$ are
$\mathcal V$ with eigenvalue $-1$ and
$\mathbb Q_\ell(1)t$ with eigenvalue $+1$.
The affine fundamental group is free, so the negative
eigenspace in its class-two quotient gives the
canonical exact sequence of local systems
\[
 0\longrightarrow \mathbb Q_\ell(1)\otimes\mathcal V
 \ \xrightarrow{\ t\otimes v\mapsto[t,v]\ }\
 (\mathfrak p/\Gamma_3\mathfrak p)^-
 \longrightarrow\mathcal V\longrightarrow0.
                                                        \tag{A.30.3}
\]
In particular no choice of a global Johnson splitting
is being assumed.

Its extension class has coefficients
$\operatorname{End}(\mathcal V)(1)$. Polarization
identifies this with $\mathcal V\otimes\mathcal V$.
Apply the alternating projection and then divide by
the polarization line to obtain a canonical class
\[
 a_\ell\in H^1_{\mathrm{cts}}
       (\mathcal H,\mathcal W_\ell),\qquad
 \mathcal W_\ell=
       \Lambda^2\mathcal V/\mathbb Q_\ell\Omega .
                                                        \tag{A.30.4}
\]
Here $\Omega$ denotes the polarization tensor, so the
displayed line is intrinsically a copy of
$\mathbb Q_\ell(1)$. We normalize the alternating
identification so that a self-adjoint symplectic
projection $A$ is sent to $\Omega_A=AJ$, in the tensor
notation of A.29.

On any subgroup acting trivially on homology, the
restriction of (A.30.3) is represented by the literal
first-Johnson homomorphism
\[
                  v\longmapsto[t,A_\gamma v].
                                                        \tag{A.30.5}
\]
Thus, when the peripheral relation makes $AJ$ skew,
the restriction of $a_\ell$ is exactly the class of
$\Omega_{A_\gamma}$ modulo $\Omega$.

\paragraph{2. The universal regulator and its precise comparison target.}
Let $\mathcal J$ be the relative Jacobian. The
Collino cycle is represented by
\[
 Z_{Q,R}=(\iota_Q(\mathcal C),f)
                  +(\iota^R(\mathcal C),f),\qquad
 \iota_Q(x)=[x-Q],\quad \iota^R(x)=[R-x].
                                                        \tag{A.30.6}
\]
The divisors cancel: they are $2[0]-2[R-Q]$ and its
negative. This defines a relative, and hence a total,
higher-Chow cycle of codimension $g$ and degree one.
The normalization $f(P)=1$ fixes its decomposable
ambiguity. The alternative second embedding $x-R$
gives the same cycle here, because the hyperelliptic
involution sends $x-R$ to $R-x$ and fixes $f$.
Rescaling $f$ by a function on the base adds a
decomposable class with coefficient twice the
polarization line, so it does not change the
primitive class below. Thus this normalization
also agrees with the primitive normal function
used in the quoted monodromy calculation.

Take the rational continuous \'etale cycle class in
$H^{2g-1}(\mathcal J,\mathbb Q_\ell(g))$. There is a
canonical way to extract its Leray degree-one part:
fix an integer $n>1$ and take the $n^{2g-2}$
eigenspace of $[n]^*$. On
$R^j\pi_*\mathbb Q_\ell(g)$ the eigenvalue is $n^j$,
so in total degree $2g-1$ this eigenspace identifies
canonically with
\[
 H^1\bigl(\mathcal H,
          R^{2g-2}\pi_*\mathbb Q_\ell(g)\bigr)
       =H^1(\mathcal H,\Lambda^2\mathcal V).
\]
Distinct eigenvalues split the finite Leray
filtration and kill its possible differentials.
Thus neither a noncanonical edge map nor a
choice of a lift through that filtration is involved.
Projecting away from $\Omega$ gives
\[
                    r_\ell(Z_{Q,R})
                      \in H^1(\mathcal H,\mathcal W_\ell).
                                                        \tag{A.30.7}
\]

We use the usual comparison compatibility of higher
cycle classes, rather than an injectivity conjecture
for realizations. In degree one this compatibility
can also be read directly from the construction:
a symbol $(Y,u)$ is the Gysin image of the Kummer
class of $u$ on the complement of its divisor;
the boundary is the divisor of $u$. The cancellation
in (A.30.6) closes the corresponding total cycle.
Over $\mathbb C$, the comparison sends the Kummer
class to the winding-number class of $u$, and
commutes with the Gysin and localization maps.
It therefore sends the \'etale class of (A.30.6)
to its Betti cycle class. The latter is the
underlying topological class of the regulator
normal function; its Leray degree-one class is
the monodromy of that normal function.
The $[n]^*$ projector commutes with all these
comparisons and with specialization.

For a published formulation of the same comparison
by universal Chern classes, see
\cite[Section 2, especially Lemma 2.2]{fghAsakuraSaito}.
That lemma is displayed there in codimension two;
its proof uses the universal Chern-class compatibility
in every degree. The degree-one symbol construction
just given supplies the required codimension-$g$
case. The passage from a universal higher-Chow
class to arithmetic by geometric comparison and
vanishing primitive invariants is also carried out
in \cite[Claim 3.3 and the following comparison on
pp.~171--172]{fghAsakuraSaito}. The analogous
fundamental-group comparison for the Ceresa class
is \cite[Corollary 6.7 and Theorem 8.1]{fghHainMatsumoto}.

\paragraph{3. The geometric comparison is a direct Dehn-twist calculation.}
Embed the base field in $\mathbb C$. The geometric
mapping class group of $\mathcal H$ is the
hyperelliptic group with all Weierstrass points
fixed; write $T\mathcal H$ for its Torelli subgroup.
Both the Betti version of (A.30.4) and the
monodromy class of (A.30.7) restrict to ordinary
homomorphisms on $T\mathcal H$.

Let $d$ be a symmetric separating simple closed
curve, with positive Dehn twist $T_d$. If $d$
does not separate $Q$ from $R$, its class in
$H_1(C\setminus\{Q,R\})$ is zero. The twist
conjugates the generators on one side by a
loop with that homology class and fixes those
on the other side, so its ordinary first
Johnson map is zero.

Suppose instead that $d$ separates $Q$ from $R$.
Write $V_Q,V_R$ for the proper homology of the
two sides and $\Omega_Q,\Omega_R$ for their
symplectic tensors. The first homology class
of the boundary of either side is $t$ or $-t$.
The same based twist formula now gives
$A_{T_d}=\pm\operatorname{pr}_{V_Q}$ or
$\mp\operatorname{pr}_{V_R}$, according to
which side contains $P$ and the fixed orientation
convention. These matrices differ, after using
one fixed convention, by a scalar identity.
Consequently their primitive tensor is one
fixed sign times $\Omega_R\bmod\Omega$.
This also verifies independence from the
position of the ordinary base point.

The actual Collino-regulator monodromy is
$4\Omega_R\bmod\Omega$ in this case, and zero
in the preceding case, by
\cite[Theorem 8.7]{fghLuoWatanabe}.
This is a statement about the regulator
normal function itself. It does not use an
identification of our open class-two extension
with a compact-curve Johnson projection.
Theorem A of
\cite{fghBrendleMargalitPutman}
says that symmetric separating twists generate
the hyperelliptic Torelli group. Hence the two
geometric classes restrict to homomorphisms
satisfying
\[
                    r_{\mathrm B}|_{T\mathcal H}
                         =\eta\,4a_{\mathrm B}|_{T\mathcal H},
                  \qquad \eta\in\{1,-1\},
                                                        \tag{A.30.8}
\]
where $\eta$ merely records a common choice of
the twist and boundary orientations.

This equality holds in $H^1$ of the full geometric
group as well. Indeed its restriction to Torelli
is injective for the coefficient representation
$\Lambda^2V/\mathbb Q\Omega$:
\cite[Theorem 5.1 and Section 5.3]{fghLuoWatanabe}
identifies this $H^1$ with the equivariant
homomorphisms out of the abelianization of
the relative unipotent radical, and proves
that the Torelli image is Zariski dense in
that radical. An equivariant homomorphism
zero on Torelli is therefore zero.
After tensoring with $\mathbb Q_\ell$, the
Betti--\'etale comparison above proves
(A.30.8) for the geometric \'etale classes.
More explicitly, the topological geometric
fundamental group has dense image in its
\'etale profinite completion. Restriction
of continuous $H^1$ to this dense subgroup
is injective: if a continuous cocycle is
the coboundary of a vector on the dense
subgroup, subtract that coboundary and
use continuity to make it zero everywhere.
Thus no separate comparison assertion for
arbitrary $\mathbb Q_\ell$ local systems
is needed at this step.

Finally the geometric homology monodromy contains
$\operatorname{Sp}_{2g}(\mathbb Z)[2]$
and is Zariski dense in $\operatorname{Sp}_{2g}$;
see \cite[Theorem 3.6 and Section 4.1]{fghLuoWatanabe}.
The only invariant in $\Lambda^2V$ is the
polarization line, so
$H^0(\mathcal H_{\overline{\mathbb Q}},
\mathcal W_\ell)=0$. The first
Hochschild--Serre terms now make
\[
 H^1(\mathcal H,\mathcal W_\ell)
       \longrightarrow
 H^1(\mathcal H_{\overline{\mathbb Q}},\mathcal W_\ell)
                                                        \tag{A.30.9}
\]
injective. We obtain the unconditional arithmetic
identity
\[
                       r_\ell(Z_{Q,R})=\eta\,4a_\ell.
                                                        \tag{A.30.10}
\]
No assertion that two arbitrary Hodge extensions
have equal \'etale realizations has been used.
In particular the corrected augmentation-space
choices in Colombo's comparison, discussed in
\cite[Introduction and Theorem 4.14]{fghSarkarSreekantan},
are not premises of this proof.

\paragraph{4. Specialization and the tangential base point.}
Specialize to $C:v^2=w^7-w$, $Q=P_0$, $R=P_\infty$,
and $P=(1,0)$. Then the normalized function is
exactly $f=w$. Over $L$ all Weierstrass points
are rational, so this is a point of $\mathcal H$.

For completeness, the primitive extension class
is unchanged when it is compared with the
tangential packet action. Rationally the
hyperelliptic involution can be made into an
actual commuting involution at a tangential
point at $Q$: its derivative is $-1$, and the
Kummer class of $-1$ vanishes in
$H^1(G_L,\mathbb Q_\ell(1))$. The resulting
invariant rational path from a tangent to its
negative has square zero as a loop, because
$\mathbb Q_\ell(1)^{G_L}=0$. Thus the corrected
involution has order two. Twisting a packet
by cusp inertia commutes with this involution.

Choose an involution-equivariant rational
unipotent path from this tangential point to
$P$. Such a path exists by successive
averaging in the central quotients of its
unipotent path torsor: the order-two group
has vanishing positive cohomology on
$\mathbb Q_\ell$-vector spaces. The two
Galois actions then differ by an inner
cocycle invariant under the involution.
Its linear part lies in the positive
eigenspace $\mathbb Q_\ell t$. Its effect
on (A.30.5) is therefore a scalar
$\operatorname{ad}(t)$, which disappears
in $\mathcal W_\ell$.
The same argument applies to every induced
original packet after the permitted common
finite descent. At $\ell=3$ the integral
packet change still has the divisibility
described in A.24; the present comparison
is rational and introduces no integral
division assertion.

It follows that (A.30.10), restricted to
any permitted subgroup of $\Delta$, is
the primitive tensor of the actual matrices
$A_\delta$ in A.29, multiplied by the
fixed nonzero rational factor $\eta4$.

\paragraph{5. A degree-four elliptic projection and its exact norm.}
Put
\[
 x_0=w-w^{-1},\quad
 x_1=\zeta w-\zeta^2w^{-1},\quad
 x_2=\zeta^2w-\zeta w^{-1},\qquad
 u=v/w^2 .
\]
Then $x_0+x_1+x_2=0$ and
$u^2=x_j^3+3x_j$ for $j=0,1,2$.
Hence $q_0=(x_0,u)$ and $q_1=(x_1,u)$
are degree-two maps to $F$.
With $[i](X,Y)=(-X,iY)$, define
$q=q_0+[i]q_1$. The addition formula gives
\[
 \begin{split}
 X_q&=\frac{-2i(w^3-w^{-3})}{x_2^2}-x_0+x_1,\\
 Y_q&=u\left(\frac{(1-i)(x_0-X_q)}{x_0+x_1}-1\right).
 \end{split}                                             \tag{A.30.11}
\]
After cancellation, $X_q=N(w)/D(w)$ with
$\deg N=4$, $\deg D=3$, $D(0)=0$, and
$D$ a nonzero constant times
$w(w^2-\zeta^2)$. The numerator and denominator
are relatively prime. Thus $q$ has degree
four and
\[
                       q^{-1}(O)
                       =\{P_0,P_\infty,P_\zeta,P_{-\zeta}\}
                                                        \tag{A.30.12}
\]
with ramification index one at all four points.
The poles of $X_q$ have order two on the
hyperelliptic curve, including at the two
finite Weierstrass points $\pm\zeta$.

For a generic target point $(X,Y)$, $Y_q=vR(w)$
recovers $v$ uniquely from $w$. The four conjugates
of $w$ are therefore exactly the roots of
$N(w)-XD(w)$. Their product is
$N(0)/\operatorname{lc}(N)$, independent of $X$.
Equivalently, the coefficients of $w$ at infinity
and $w^{-1}$ at zero in $X_q$ are
\[
          a_\infty=\zeta-1-2i\zeta^2,\qquad
          a_0=1-\zeta^2+2i\zeta .
\]
Since $D$ is proportional to $w(w^2-\zeta^2)$,
\[
 \operatorname{Norm}_{L(C)/L(F)}(w)
       =-\zeta^2\,\frac{a_0}{a_\infty}
       =\zeta(2+\sqrt3)^{-2}.
                                                        \tag{A.30.13}
\]
These are exact algebraic identities.
The retained program \texttt{code/a\_C3\_collino\_unit.py}
verifies (A.30.11), the target equation, cancellation,
degrees, all four poles, and (A.30.13) over
$\mathbb Q(i,\sqrt3)(w)$.

Both $Q$ and $R$ map to $O$ under $q$.
The two symbols in (A.30.6) consequently push
forward to the same constant norm, so their
product is $\zeta^2\epsilon^{-4}$, as claimed
in (A.30.2). For comparison, the equation
$w^2-Xw-1=0$ for $q_0$ gives
$\operatorname{Norm}_{q_0}(w)=-1$ and
$(q_0)_*Z_{Q,R}=1$.

\paragraph{6. Removing the scalar part and proving nonvanishing on $\Delta$.}
A map $C\to F$ sending $Q$ to $O$ induces a
homomorphism $J(C)\to F$, and therefore a
functional
\[
       \Lambda^2H_1(C,\mathbb Q_\ell)
                  \longrightarrow\mathbb Q_\ell(1).
\]
It sends $\Omega_C$ to
$\deg(q)\Omega_F$. This follows either from
$q_*q^*=\deg(q)$ on $H^1(F)$ and the
intersection pairing, or from the adjoint
description of the induced Jacobian
homomorphism. Consequently
\[
                        \Phi=q_*-2(q_0)_*
                                                        \tag{A.30.14}
\]
kills $\Omega_C$, because $4-2\cdot2=0$.
It is an algebraic functional on the
primitive quotient, defined over $L$.
Functoriality of higher cycle classes and
(A.30.13) give
\[
             \Phi\bigl(r_\ell(Z_{Q,R})\bigr)
                    =-4\kappa_\epsilon
                    \quad\hbox{in }H^1(G_L,\mathbb Q_\ell(1)).
                                                        \tag{A.30.15}
\]
The root of unity $\zeta^2$ has zero rational
Kummer class. This calculation is why an
off-diagonal elliptic map is needed: the
individual reflection projection has zero
regulator, while the combination (A.30.14)
also removes the scalar polarization
component.

The unit $\epsilon=2+\sqrt3>1$ has infinite
order, so $\kappa_\epsilon\ne0$ over $L$:
the unit group modulo torsion is a finitely
generated free abelian group, whose
$\ell$-adic Kummer map is injective after
tensoring with $\mathbb Q_\ell$.
Moreover $L\subset K^{\mathrm{ab}}$ for
$K=\mathbb Q(i)$. Thus $\Delta\subset G_L$
and $G_L/\Delta$ is abelian. Some element
of this quotient has cyclotomic character
different from one. Sah's cocycle identity
therefore gives
\[
                   H^1(G_L/\Delta,\mathbb Q_\ell(1))=0,
\]
since multiplication by that nonzero
$\ell$-adic scalar is invertible.
Inflation--restriction makes
\[
 H^1(G_L,\mathbb Q_\ell(1))
           \longrightarrow H^1(\Delta,\mathbb Q_\ell(1))
\]
injective. The restriction of $\kappa_\epsilon$
is hence nonzero. On $\Delta$ the coefficients
are trivial, so it is an additive character;
such a character cannot vanish on an open
subgroup unless it vanishes everywhere.
This proves nonvanishing on every permitted
$\Delta_0$.

By (A.30.10), (A.30.15), and the base-point
comparison, the actual primitive first
Johnson homomorphism is nonzero on every
such subgroup, for every packet. In particular
some actual $A_\delta$ is nonzero.
Statement A.29 now applies and gives
(A.30.1). Its integral cusp-power subtraction
is already part of that statement, so no
fixed element of a rational completion is
silently lifted to the pro-$\ell$ group.
\end{proof}

